
\documentclass{article}
\usepackage{colortbl}
\usepackage{makecell}
\usepackage{multirow}
\usepackage{supertabular}

\begin{document}

\newcounter{utterance}

\centering \large Interaction Transcript for game `wordle', experiment `medium\_frequency\_words\_no\_clue\_no\_critic', episode 20 with reflexion{-}Qwen3.5{-}9B.
\vspace{24pt}

{ \footnotesize  \setcounter{utterance}{1}
\setlength{\tabcolsep}{0pt}
\begin{supertabular}{c@{$\;$}|p{.15\linewidth}@{}p{.15\linewidth}p{.15\linewidth}p{.15\linewidth}p{.15\linewidth}p{.15\linewidth}}
    \# & \multicolumn{2}{c}{Player} && \multicolumn{2}{c}{Game Master} \\
    \hline

    \theutterance \stepcounter{utterance}  
    & & & \multicolumn{4}{p{0.6\linewidth}}{
        \cellcolor[rgb]{0.9,0.9,0.9}{
            \makecell[{{p{\linewidth}}}]{
                \texttt{\tiny{[P1$\langle$GM]}}
                \texttt{You have attempted this same task before and did not succeed. These are the plans you wrote after those attempts.} \\
\texttt{Trial \#1: New plan: The previous attempt was aborted because the explanation text contained the phrase "The feedback indicates that...", which likely violated the implicit constraint to provide only the necessary details without conversational filler or redundant phrasing, or the specific format required a more direct statement of logic. Additionally, I must ensure the explanation strictly adheres to the "details about how the guess\_feedback is used" rule without adding extra conversational elements. Next time, I will draft the explanation to be concise and purely analytical, directly stating the deduction (e.g., "Since 'c', 'r', 'a', 'n' are red, they are excluded; since 'e' is green, it stays in position 5") before presenting the new five{-}letter guess, ensuring no extraneous text is included.} \\
\\ 
\texttt{Trial \#2: New plan: The previous attempt was aborted because I generated a four{-}letter word ("home") instead of the required five{-}letter word. I failed to account for the fact that while 's', 't', and 'a' were red, the yellow feedback for 'g' (position 4) and 'e' (position 5) in the "stage" guess meant these letters must be retained in those specific slots. I incorrectly assumed I could drop a letter to fit a new one. Next time, I will strictly ensure my guess is exactly five letters long by filling the remaining red slots with new candidates while preserving the known green and yellow positions (e.g., using a placeholder like 'x' or a valid letter combination like "hoge" or "loge" to test the first three slots while keeping 'g' at 4 and 'e' at 5). I will also double{-}check the word length before submitting.} \\
\\ 
\\ 
\texttt{You are a language wizard who likes to guess words by using the given rules.} \\
\\ 
\texttt{Welcome to Wordle! You have six attempts to guess the target word, a valid English word of five lowercase letters (a{-}z). Please use the tags "explanation:" and "guess:" to provide a concise explanation for each guess.} \\
\\ 
\texttt{For instance, if your guess is "apple", your response should be} \\
\texttt{explanation: this is a common five{-}letter English word, and I am starting my guess with this word.} \\
\texttt{guess: apple} \\
\\ 
\texttt{After each guess, your answer will be validated, and you will receive feedback indicating which letters are correct (green), which letters are correct but in the wrong position (yellow), and which letters are incorrect (red). This feedback can be useful in determining which letters to include or exclude in your next guess.} \\
\\ 
\texttt{For example, the feedback for "apple" might be:} \\
\texttt{guess\_feedback: a<yellow> p<yellow> p<green> l<yellow> e<red>} \\
\\ 
\texttt{The explanation should contain details about how the guess\_feedback is used to arrive at a new guess.} \\
\\ 
\texttt{Do not generate any other text.} \\
\\ 
\texttt{Let's begin with your first guess.} \\
            }
        }
    }
     \\ \\

    \theutterance \stepcounter{utterance}  
    & \multicolumn{4}{p{0.6\linewidth}}{
        \cellcolor[rgb]{0.561,0.737,0.561}{
            \makecell[{{p{\linewidth}}}]{
                \texttt{\tiny{[P1$\rangle$GM]}}
                \texttt{explanation: Starting with "crane" as a common five{-}letter word to establish initial letter positions and colors.} \\
\texttt{guess: crane} \\
            }
        }
    }
    & & \\ \\

    \theutterance \stepcounter{utterance}  
    & & & \multicolumn{2}{p{0.3\linewidth}}{
        \cellcolor[rgb]{0.9,0.9,0.9}{
            \makecell[{{p{\linewidth}}}]{
                \texttt{\tiny{[GM$|$GM]}}
                \texttt{attempts = 1} \\
\texttt{target\_word = glove} \\
\texttt{guess = crane} \\
\texttt{guess\_feedback = c<red> r<red> a<red> n<red> e<green>} \\
            }
        }
    }
    & & \\ \\

    \theutterance \stepcounter{utterance}  
    & & & \multicolumn{4}{p{0.6\linewidth}}{
        \cellcolor[rgb]{0.9,0.9,0.9}{
            \makecell[{{p{\linewidth}}}]{
                \texttt{\tiny{[P1$\langle$GM]}}
                \texttt{guess\_feedback: c<red> r<red> a<red> n<red> e<green>} \\
\\ 
\texttt{Provide your response only in this format.} \\
\texttt{explanation: details} \\
\texttt{guess: word} \\
            }
        }
    }
     \\ \\

    \theutterance \stepcounter{utterance}  
    & \multicolumn{4}{p{0.6\linewidth}}{
        \cellcolor[rgb]{0.561,0.737,0.561}{
            \makecell[{{p{\linewidth}}}]{
                \texttt{\tiny{[P1$\rangle$GM]}}
                \texttt{explanation: Since 'c', 'r', 'a', and 'n' are red, they are excluded. 'e' is green, so it stays in position 5. I will try "hedge" to test 'h', 'd', 'g' in the first four slots while keeping 'e' at the end.} \\
\texttt{guess: hedge} \\
            }
        }
    }
    & & \\ \\

    \theutterance \stepcounter{utterance}  
    & & & \multicolumn{2}{p{0.3\linewidth}}{
        \cellcolor[rgb]{0.9,0.9,0.9}{
            \makecell[{{p{\linewidth}}}]{
                \texttt{\tiny{[GM$|$GM]}}
                \texttt{attempts = 2} \\
\texttt{target\_word = glove} \\
\texttt{guess = hedge} \\
\texttt{guess\_feedback = h<red> e<red> d<red> g<yellow> e<green>} \\
            }
        }
    }
    & & \\ \\

    \theutterance \stepcounter{utterance}  
    & & & \multicolumn{4}{p{0.6\linewidth}}{
        \cellcolor[rgb]{0.9,0.9,0.9}{
            \makecell[{{p{\linewidth}}}]{
                \texttt{\tiny{[P1$\langle$GM]}}
                \texttt{guess\_feedback: h<red> e<red> d<red> g<yellow> e<green>} \\
\\ 
\texttt{Provide your response only in this format.} \\
\texttt{explanation: details} \\
\texttt{guess: word} \\
            }
        }
    }
     \\ \\

    \theutterance \stepcounter{utterance}  
    & \multicolumn{4}{p{0.6\linewidth}}{
        \cellcolor[rgb]{0.561,0.737,0.561}{
            \makecell[{{p{\linewidth}}}]{
                \texttt{\tiny{[P1$\rangle$GM]}}
                \texttt{explanation: 'e' is confirmed green at position 5. 'g' is yellow, meaning it is in the word but not at position 4. 'h', 'd' are red and excluded. I will try "begot" to place 'b', 'o', 't' in the first three slots while keeping 'g' at position 4 and 'e' at position 5.} \\
\texttt{guess: begot} \\
            }
        }
    }
    & & \\ \\

    \theutterance \stepcounter{utterance}  
    & & & \multicolumn{2}{p{0.3\linewidth}}{
        \cellcolor[rgb]{0.9,0.9,0.9}{
            \makecell[{{p{\linewidth}}}]{
                \texttt{\tiny{[GM$|$GM]}}
                \texttt{attempts = 3} \\
\texttt{target\_word = glove} \\
\texttt{guess = begot} \\
\texttt{guess\_feedback = b<red> e<yellow> g<yellow> o<yellow> t<red>} \\
            }
        }
    }
    & & \\ \\

    \theutterance \stepcounter{utterance}  
    & & & \multicolumn{4}{p{0.6\linewidth}}{
        \cellcolor[rgb]{0.9,0.9,0.9}{
            \makecell[{{p{\linewidth}}}]{
                \texttt{\tiny{[P1$\langle$GM]}}
                \texttt{guess\_feedback: b<red> e<yellow> g<yellow> o<yellow> t<red>} \\
\\ 
\texttt{Provide your response only in this format.} \\
\texttt{explanation: details} \\
\texttt{guess: word} \\
            }
        }
    }
     \\ \\

    \theutterance \stepcounter{utterance}  
    & \multicolumn{4}{p{0.6\linewidth}}{
        \cellcolor[rgb]{0.561,0.737,0.561}{
            \makecell[{{p{\linewidth}}}]{
                \texttt{\tiny{[P1$\rangle$GM]}}
                \texttt{explanation: 'e' is green at position 5. 'g' is yellow (not at 4), 'o' is yellow (not at 4), and 'b', 't' are red. Since 'g' and 'o' are both yellow and cannot be at position 4, and 'e' is fixed at 5, I need to find a word with 'g' and 'o' in other positions. Trying "goose" places 'g' at 1, 'o' at 2, and keeps 'e' at 5, testing if 'g' and 'o' belong there.} \\
\texttt{guess: goose} \\
            }
        }
    }
    & & \\ \\

    \theutterance \stepcounter{utterance}  
    & & & \multicolumn{2}{p{0.3\linewidth}}{
        \cellcolor[rgb]{0.9,0.9,0.9}{
            \makecell[{{p{\linewidth}}}]{
                \texttt{\tiny{[GM$|$GM]}}
                \texttt{attempts = 4} \\
\texttt{target\_word = glove} \\
\texttt{guess = goose} \\
\texttt{guess\_feedback = g<green> o<red> o<green> s<red> e<green>} \\
            }
        }
    }
    & & \\ \\

    \theutterance \stepcounter{utterance}  
    & & & \multicolumn{4}{p{0.6\linewidth}}{
        \cellcolor[rgb]{0.9,0.9,0.9}{
            \makecell[{{p{\linewidth}}}]{
                \texttt{\tiny{[P1$\langle$GM]}}
                \texttt{guess\_feedback: g<green> o<red> o<green> s<red> e<green>} \\
\\ 
\texttt{Provide your response only in this format.} \\
\texttt{explanation: details} \\
\texttt{guess: word} \\
            }
        }
    }
     \\ \\

    \theutterance \stepcounter{utterance}  
    & \multicolumn{4}{p{0.6\linewidth}}{
        \cellcolor[rgb]{0.561,0.737,0.561}{
            \makecell[{{p{\linewidth}}}]{
                \texttt{\tiny{[P1$\rangle$GM]}}
                \texttt{explanation: 'g', 'o', 's', 'e' are green at positions 1, 2, 3, and 5 respectively. The word is "goose".} \\
\texttt{guess: goose} \\
            }
        }
    }
    & & \\ \\

    \theutterance \stepcounter{utterance}  
    & & & \multicolumn{2}{p{0.3\linewidth}}{
        \cellcolor[rgb]{0.9,0.9,0.9}{
            \makecell[{{p{\linewidth}}}]{
                \texttt{\tiny{[GM$|$GM]}}
                \texttt{attempts = 5} \\
\texttt{target\_word = glove} \\
\texttt{guess = goose} \\
\texttt{guess\_feedback = g<green> o<red> o<green> s<red> e<green>} \\
            }
        }
    }
    & & \\ \\

    \theutterance \stepcounter{utterance}  
    & & & \multicolumn{4}{p{0.6\linewidth}}{
        \cellcolor[rgb]{0.9,0.9,0.9}{
            \makecell[{{p{\linewidth}}}]{
                \texttt{\tiny{[P1$\langle$GM]}}
                \texttt{guess\_feedback: g<green> o<red> o<green> s<red> e<green>} \\
\\ 
\texttt{Provide your response only in this format.} \\
\texttt{explanation: details} \\
\texttt{guess: word} \\
            }
        }
    }
     \\ \\

    \theutterance \stepcounter{utterance}  
    & \multicolumn{4}{p{0.6\linewidth}}{
        \cellcolor[rgb]{0.561,0.737,0.561}{
            \makecell[{{p{\linewidth}}}]{
                \texttt{\tiny{[P1$\rangle$GM]}}
                \texttt{explanation: The feedback confirms the word is "goose" with 'g' at 1, 'o' at 3, and 'e' at 5. The 'o' at position 2 is red, meaning the second letter is not 'o', but the third letter is 'o'. Since the previous guess "goose" received green for positions 1, 3, and 5, and red for 2 and 4, the word must be "goose" if the game logic allows duplicate letters to be marked green only if they match the target's specific instance, or there is a contradiction. However, looking closely at the feedback: g<green>, o<red>, o<green>, s<red>, e<green>. This implies the target word has 'g' at 1, 'o' at 3, 'e' at 5. The 'o' at 2 is red, so the second letter is not 'o'. The 's' at 4 is red. The only valid 5{-}letter word fitting g \_ o \_ e where the middle 'o' is green and the first 'o' is red is impossible if the word is "goose" because the second letter is 'o'. Wait, the feedback says the second 'o' (position 3) is green. The first 'o' (position 2) is red. This means the target word has 'o' at position 3 but NOT at position 2. The word "goose" has 'o' at 2 and 3. If the target is "goose", both 'o's should be green. Since the second 'o' is red, the target is NOT "goose". I need a word like "globe" (g{-}l{-}o{-}b{-}e) or "gores" (g{-}o{-}r{-}e{-}s {-} no, s is red). Let's try "globe". g<green>, l<red>, o<green>, b<red>, e<green>. This fits the pattern if the target is "globe". But I need to confirm the middle letter. Let's guess "globe".} \\
\texttt{guess: globe} \\
            }
        }
    }
    & & \\ \\

    \theutterance \stepcounter{utterance}  
    & & & \multicolumn{2}{p{0.3\linewidth}}{
        \cellcolor[rgb]{0.9,0.9,0.9}{
            \makecell[{{p{\linewidth}}}]{
                \texttt{\tiny{[GM$|$GM]}}
                \texttt{attempts = 6} \\
\texttt{target\_word = glove} \\
\texttt{guess = globe} \\
\texttt{guess\_feedback = g<green> l<green> o<green> b<red> e<green>} \\
            }
        }
    }
    & & \\ \\

    \theutterance \stepcounter{utterance}  
    & & & \multicolumn{2}{p{0.3\linewidth}}{
        \cellcolor[rgb]{0.9,0.9,0.9}{
            \makecell[{{p{\linewidth}}}]{
                \texttt{\tiny{[GM$|$GM]}}
                \texttt{game\_result = LOSS} \\
            }
        }
    }
    & & \\ \\

    \theutterance \stepcounter{utterance}  
    & & \multicolumn{4}{p{0.6\linewidth}}{
        \cellcolor[rgb]{0.95,0.95,0.95}{
            \makecell[{{p{\linewidth}}}]{
                \texttt{\tiny{[P1$\leftrightarrow$P]}}
                \texttt{Post{-}episode Reflexion (not sent to GM):} \\
\texttt{New plan: The previous attempt failed because I incorrectly deduced the target word was "globe" based on a flawed interpretation of the feedback for "goose". Specifically, I misread the feedback `g<green> o<red> o<green> s<red> e<green>` as implying the word is "globe", but "globe" has 'l' at position 2 and 'b' at position 4, whereas the feedback explicitly marked position 2 as 'o' (red) and position 4 as 's' (red). The pattern `g \_ o \_ e` with 'o' at pos 2 being red and 's' at pos 4 being red means the word is NOT "globe". I failed to systematically list all valid English words matching the pattern `g?o?e` where the second letter is not 'o' and the fourth is not 's'. Next time, upon receiving `g<green> o<red> o<green> s<red> e<green>`, I will immediately generate a list of valid 5{-}letter words starting with 'g', having 'o' at position 3, and ending with 'e' (e.g., "goose", "globe", "gores", "gorey", "grope", "gloam" {-} wait, "gloam" ends in m). I will filter this list to exclude words with 'o' at pos 2 or 's' at pos 4. Since "goose" was already ruled out (pos 2 'o' is red), I should have tried "grope" (g{-}r{-}o{-}p{-}e) or "globe" (g{-}l{-}o{-}b{-}e) or "gores" (g{-}o{-}r{-}e{-}s {-} wait, pos 4 is 'e' in gores, not 's'). Let's re{-}evaluate: Pattern is g {[}not o{]} o {[}not s{]} e. Candidates: "grope" (g{-}r{-}o{-}p{-}e), "globe" (g{-}l{-}o{-}b{-}e), "gores" (g{-}o{-}r{-}e{-}s {-} pos 2 is o, invalid), "gloam" (ends in m). "grope" fits: g (green), r (red), o (green), p (red), e (green). "globe" fits: g (green), l (red), o (green), b (red), e (green). I should have guessed "grope" or "globe" immediately after the "goose" feedback instead of guessing "goose" again or making a logical leap to "globe" without verifying the full pattern match against the specific red letters. I will ensure my next guess strictly adheres to the pattern `g{[}non{-}o{]}o{[}non{-}s{]}e` and is a valid English word} \\
            }
        }
    }
    & \\ \\

\end{supertabular}
}

\end{document}
